\section{Proofs}\label{app:proofs}

\subsection{Setting and notation}

The proofs use the following, more general, statement of the setting. Rows $i=1,\dots,n$ have outcomes
$y_i=\mu_i+\varepsilon_i$ with $\mu_i=x_i'\beta$; the regressors and $\beta$ are non-random. The estimator is
$\hat\theta=\sum_i\sum_\ell C_{i\ell}y_iy_\ell$, with $m=C\mu$, $\hat m=Cy$ and $V=\V[\hat\theta]$. $U$ denotes the
set of rows of a stayer, $\mathcal{U}$ the collection of these sets (one per stayer), $|U|\ge2$ the number of rows
of $U$ and $d_U=|U|-1$; $R$ is the set of movers' rows ($\cR$ in the text).

\begin{enumerate}[label=\textbf{(S\arabic*)},leftmargin=3.3em,itemsep=3pt]
\item \textbf{Independent errors.} $\varepsilon_1,\dots,\varepsilon_n$ are independent, with $\E\varepsilon_i=0$
and $0<\sigma_i^2=\E\varepsilon_i^2<\infty$.
\item \textbf{Exact weights.} $C$ is a non-random symmetric matrix with $C_{ii}=0$ for all $i$, and for every
$U\in\mathcal{U}$ there is a number $c_U$ with $C_{i\ell}=c_U$ for all $i\ne\ell$ in $U$. (Lemma~\ref{lem:stayer}:
the weights \eqref{eq:ustat} have this property.)
\item \textbf{Stayers.} The sets $U\in\mathcal{U}$ are pairwise disjoint and are determined by the regressors.
Within a set the means are equal, $\mu_i=\mu_U$, and so are the variances, $\sigma_i^2=\tau_U$. With
$\bar y_U$ the mean of the outcomes of $U$,
\[
q_U=\frac{1}{d_U}\sum_{i\in U}(y_i-\bar y_U)^2,
\]
and every row of $U$ receives $\st_i^2=q_U$. ($q_U$ is $q_g$ of \eqref{eq:q} and $\tau_U$ is $\sigma_g^2$.)
\item \textbf{Stayers' errors.} For every row $i$ of a set, $\E\varepsilon_i^3=0$,
$\E\varepsilon_i^4=\kappa\,\sigma_i^4$ and $\E\varepsilon_i^8\le\bar c\,\sigma_i^8$. The text takes $\kappa=3$;
the proofs allow any $\kappa\ge1$.
\item \textbf{Movers' rows.} For $i\in R$, $\E\varepsilon_i^4\le\bar\kappa\,\sigma_i^4$ for a constant
$\bar\kappa\ge\kappa$. These rows receive $\st_i^2=\sigma_i^2+\Delta_i$ with non-random numbers $\Delta_i$.
\item $V>0$.
\end{enumerate}
S1 to S3 are A1 and A2 of the text, S4 is A3 with the moment bound of Proposition~\ref{prop:conc}, and nothing
is asked of the third moments or the kurtosis of movers' errors.

\paragraph{The estimator.} For a stated kurtosis $\kappa$ let
\begin{equation}\label{eq:nuh}
\nu_U=\frac{2}{|U|-1}+\frac{\kappa-3}{|U|},\qquad
\chi_U=c_U^2\big[4\,(2-|U|)+2\,(\kappa-3)\,(|U|-1)\big],\qquad \bnu=2+|\kappa-3| ,
\end{equation}
and
\begin{equation}\label{eq:VhatA}
\hat V=4\sum_i\hat m_i^{\,2}\,\st_i^2-2\sum_i\sum_\ell C_{i\ell}^2\,\st_i^2\,\st_\ell^2
-\sum_{U\in\mathcal{U}}\chi_U\,\frac{q_U^2}{1+\nu_U}.
\end{equation}
For $\kappa=3$, $\chi_U/(1+\nu_U)=-4c_U^2(|U|-2)(|U|-1)/(|U|+1)$, and \eqref{eq:VhatA} is
$\hat\V[\hat\theta]$ of \eqref{eq:Vhat}. $b(\Delta)$ is defined in \eqref{eq:b}.

\paragraph{Standardised notation.} Let $u_i=\varepsilon_i/\sigma_i$, $a_i=\sigma_im_i$ and
$K_{i\ell}=\sigma_iC_{i\ell}\sigma_\ell$, so that $K$ is symmetric with zero diagonal, and let $w=a+Ku$, so that
$\sigma_i\hat m_i=w_i$. For index sets $I,J$, $K_{IJ}$ is the block of $K$ with rows in $I$ and columns in $J$,
$K_{I\cdot}$ the rows in $I$, and $u_I$, $a_I$, $w_I$ the subvectors. $\|\cdot\|$ is the Euclidean norm of a vector
and the Frobenius norm of a matrix; $\|\cdot\|_\op$ is the operator norm. We use $\|XY\|\le\|X\|_\op\|Y\|$,
$\|X\|_\op\le\|X\|$ and $|\tr(XY)|\le\|X\|\,\|Y\|$ freely.

For $U\in\mathcal{U}$ let $\Pi_U=I-\mathbf{1}\mathbf{1}'/|U|$, the projection that removes the mean of the set, and
$\bar Q_U=\Pi_U/d_U$. Since the means are equal within $U$ and $\varepsilon_i=\sqrt{\tau_U}\,u_i$ on $U$,
\begin{equation}\label{eq:xi}
\xi_U:=\frac{q_U}{\tau_U}-1=u_U'\,\bar Q_U\,u_U-1,\qquad \st_i^2=\sigma_i^2\,(1+\xi_U)\ \ (i\in U).
\end{equation}
$\xi_U$ depends on $u_U$ only, so the $\xi_U$ are independent across sets, and $\xi_U$ is independent of the errors
outside $U$. Put $\xi_i=\xi_U$ for $i\in U$ and $\xi_i=0$ for $i\in R$, and let
\begin{equation}\label{eq:G}
G_U=2\,\bar Q_U+(\kappa-3)\diag\big(\bar Q_U\big)=2\,\bar Q_U+\frac{\kappa-3}{|U|}\,I .
\end{equation}
Note that $\tr(\bar Q_U^2)=1/d_U$ and that the diagonal entries of $\bar Q_U$ equal $1/|U|$, so that
$\nu_U=2\tr(\bar Q_U^2)+(\kappa-3)\sum_{j\in U}\bar Q_{U,jj}^{\,2}$.

For the movers' rows let $\vartheta_i=1+\Delta_i/\sigma_i^2$ ($i\in R$), $\vartheta_i=1$ for rows in sets,
$\Theta=\diag(\vartheta_i)$ and $\bth=\max(1,\max_{i\in R}|\vartheta_i|)$. Then
$\st_i^2=\sigma_i^2\,\vartheta_i\,(1+\xi_i)$ for every row, and \eqref{eq:VhatA} reads
\begin{equation}\label{eq:pieces}
\begin{aligned}
\hat V&=P_1+P_2+P_3+P_4, &
P_1&=4\,w'\Theta w, \qquad P_2=4\sum_{U\in\mathcal{U}}\xi_U\,\|w_U\|^2,\\
P_3&=-2\sum_i\sum_\ell K_{i\ell}^2\,\vartheta_i\vartheta_\ell\,(1+\xi_i)(1+\xi_\ell), &
P_4&=-\sum_{U\in\mathcal{U}}\tau_U^2\,\chi_U\,\frac{(1+\xi_U)^2}{1+\nu_U}.
\end{aligned}
\end{equation}
(For the first term, $\hat m_i^{\,2}\st_i^2=w_i^2\vartheta_i(1+\xi_i)$, and $\vartheta_i=1$ where $\xi_i\ne0$.)
By Lemma~\ref{lem:forms} below, $V=4\|a\|^2+2\|K\|^2$, which is \eqref{eq:V}. Let $k_{UU'}=\|K_{UU'}\|^2$,
$k_U=\|K_{U\cdot}\|^2=\sum_{i\in U}\sum_\ell K_{i\ell}^2$, $\rho_U=\|a_U\|^2+k_U$, and let
$\lambda=\lambda_{\max}=\|K\|_\op$. The share of Proposition~\ref{prop:conc} is $\phi_U=(4\|a_U\|^2+2k_U)/V$. Since
the sets are disjoint,
\begin{equation}\label{eq:facts}
2\rho_U\le\phi_UV,\qquad \sum_{U}\rho_U\le\|a\|^2+\|K\|^2\le\tfrac12V,\qquad
\sum_{U'\in\mathcal{U}}k_{UU'}\le k_U,\qquad k_{UU}\le k_U\le\rho_U .
\end{equation}
All results are bounds for one design with explicit constants; along a sequence of designs the constants of
the assumptions are the same for every design.

\subsection{Lemmas}

\begin{lemma}[A stayer's weights]\label{lem:stayer}
\textup{(a)} For the weights \eqref{eq:ustat} and two rows $i\ne\ell$ of stayer $g$: $P_{i\ell}=1/T_g$,
$B_{i\ell}=A_{gg}/T_g^2$ and $C_{i\ell}=A_{gg}/(T_g(T_g-1))$; moreover $m_i=m_\ell$.
\textup{(b)} Under S2, $K_{UU}=\tau_Uc_U(\mathbf{1}\mathbf{1}'-I)$, and with $e_U:=\tr(K_{UU}^2G_U)$,
\[
4\,e_U-2\,\nu_U\,k_{UU}=\tau_U^2\,\chi_U,\qquad \tau_U^2\,|\chi_U|\le\bnu\,k_{UU} .
\]
\end{lemma}

\begin{proof}
(a) The rows of stayer $g$ have the same regressor vector $x_i=x_g$, and no other row involves the person effect
of $g$. Let $\epsilon_g$ be the unit vector of that effect. Then
$S_{xx}\epsilon_g=\sum_jx_j(x_j'\epsilon_g)=T_gx_g$, so $S_{xx}^{-1}x_i=\epsilon_g/T_g$. Hence
$P_{i\ell}=x_i'\epsilon_g/T_g=1/T_g$ and $B_{i\ell}=\epsilon_g'A\epsilon_g/T_g^2=A_{gg}/T_g^2$, also for $i=\ell$.
With $M_{i\ell}=-1/T_g$ and $M_{ii}=1-1/T_g$, \eqref{eq:ustat} gives
$C_{i\ell}=\frac{A_{gg}}{T_g^2}\big(1+\frac{1}{T_g-1}\big)=\frac{A_{gg}}{T_g(T_g-1)}$. Finally
$m_i-m_\ell=\sum_{j\ne i}C_{ij}\mu_j-\sum_{j\ne\ell}C_{\ell j}\mu_j=C_{i\ell}(\mu_\ell-\mu_i)=0$, because the rows
$i$ and $\ell$ of $C$ agree outside the columns $i$ and $\ell$ and $\mu_i=\mu_\ell$.
(b) Write $t=|U|$ and $J=\mathbf{1}\mathbf{1}'$. The first claim is S2 with $\sigma_i^2=\tau_U$ on $U$. Then
$K_{UU}^2=\tau_U^2c_U^2\big((t-2)J+I\big)$, so $k_{UU}=\tr K_{UU}^2=\tau_U^2c_U^2\,t(t-1)$, and, since $J\Pi_U=0$,
$\tr(K_{UU}^2\bar Q_U)=\tau_U^2c_U^2\tr(\Pi_U)/d_U=\tau_U^2c_U^2$. By \eqref{eq:G} and \eqref{eq:nuh},
\[
4e_U-2\nu_Uk_{UU}=\tau_U^2c_U^2\Big[8+4(\kappa-3)(t-1)-4t-2(\kappa-3)(t-1)\Big]=\tau_U^2\chi_U .
\]
For the bound, $|2-t|\le t-1$ gives $|\chi_U|\le c_U^2(t-1)(4+2|\kappa-3|)=2\bnu\,c_U^2(t-1)$, and
$c_U^2(t-1)\tau_U^2=k_{UU}/t\le k_{UU}/2$.
\end{proof}

In Lemma~\ref{lem:forms}, $u_1,\dots,u_N$ are independent with $\E u_j=0$, $\E u_j^2=1$ and
$\kappa_j:=\E u_j^4\le\bar\kappa$. Nothing is assumed about third moments. Put
$c_4=\max(3,\bar\kappa)$.

\begin{lemma}[Linear and quadratic forms]\label{lem:forms}
Let $b\in\mathbb{R}^N$, $M$ a symmetric $N\times N$ matrix, $a\in\mathbb{R}^t$ and $F$ a $t\times N$
matrix, all non-random.
\begin{enumerate}[label=\textup{(\alph*)}]
\item $\E(b'u)^2=\|b\|^2$ and $\E(b'u)^4=3\|b\|^4+\sum_j(\kappa_j-3)\,b_j^4\le c_4\|b\|^4$.
\item $\E[u'Mu]=\tr M$ and $\Var(u'Mu)=2\|M\|^2+\sum_j(\kappa_j-3)M_{jj}^2\le\max(2,\bar\kappa-1)\,\|M\|^2$.
If $M$ has zero diagonal, then $\Var(u'Mu)=2\|M\|^2$ and $\E[(b'u)(u'Mu)]=0$, and these two
statements need second moments only.
\item $\E\|a+Fu\|^4\le4c_4\big(\|a\|^2+\|F\|^2\big)^2$.
\item If $M$ has zero diagonal, $\E(u'Mu)^4\le(256+24\bar\kappa)\,c_4\,\|M\|^4$.
\item If in addition $\E u_j^8\le\bar c$ for all $j$, then $\E(u'Mu-\tr M)^4\le c_Q\,\|M\|^4$ with
$c_Q=6720\,\bar c$. If $M$ is also positive semidefinite, $\E(u'Mu)^4\le8\,(c_Q+1)\,(\tr M)^4$.
\end{enumerate}
\end{lemma}

\begin{proof}
A product of errors has expectation zero whenever some error occurs in it exactly once (independence
and mean zero). Call this the \emph{single-index rule}.

(a) The second moment is immediate. In
$\E(b'u)^4=\sum b_{j_1}b_{j_2}b_{j_3}b_{j_4}\E[u_{j_1}u_{j_2}u_{j_3}u_{j_4}]$ the single-index rule
leaves the terms with four equal indices, $\sum_jb_j^4\kappa_j$, and those with two distinct indices
each occurring twice, $3\sum_{j\ne\ell}b_j^2b_\ell^2=3\|b\|^4-3\sum_jb_j^4$. For the bound,
$\sum_jb_j^4\le\|b\|^4$, so the fourth moment is at most $3\|b\|^4$ if $\bar\kappa\le3$ and at most
$\bar\kappa\|b\|^4$ otherwise.

(b) Write $u'Mu-\tr M=D+O$ with $D=\sum_jM_{jj}(u_j^2-1)$ and $O=2\sum_{j<\ell}M_{j\ell}u_ju_\ell$.
The summands of $D$ are independent with mean zero and variances $M_{jj}^2(\kappa_j-1)$. Each summand
of $O$ has mean zero and variance $4M_{j\ell}^2$, and by the single-index rule it is uncorrelated
with every other summand of $O$, with every summand of $D$ and with $b'u$. Hence
$\Var(u'Mu)=\sum_jM_{jj}^2(\kappa_j-1)+2\sum_{j\ne\ell}M_{j\ell}^2$, which is the formula. It is
at most $2\|M\|^2$ if $\bar\kappa\le3$ and at most $2\|M\|^2+(\bar\kappa-3)\sum_jM_{jj}^2
\le(\bar\kappa-1)\|M\|^2$ otherwise. With zero diagonal, $D=0$ and no fourth moment enters.

(c) Let $g=a+Fu$ with coordinates $g_i=a_i+f_i'u$, $f_i'$ the $i$th row of $F$. By Minkowski's
inequality in $L^4$ and (a), $(\E g_i^4)^{1/4}\le|a_i|+c_4^{1/4}\|f_i\|$, so
$(\E g_i^4)^{1/2}\le2a_i^2+2c_4^{1/2}\|f_i\|^2\le2c_4^{1/2}(a_i^2+\|f_i\|^2)$. By Minkowski's inequality in
$L^2$, $(\E\|g\|^4)^{1/2}=\big(\E(\sum_ig_i^2)^2\big)^{1/2}\le\sum_i(\E g_i^4)^{1/2}
\le2c_4^{1/2}(\|a\|^2+\|F\|^2)$.

(d) Let $L_k=\sum_{i<k}M_{ik}u_i$ (with $L_1=0$), $\gamma_k=2u_kL_k$ and $S_k=\sum_{j\le k}\gamma_j$, so that
$u'Mu=S_N$. $L_k$ is a function of $u_1,\dots,u_{k-1}$, so $\E[\gamma_k\mid u_1,\dots,u_{k-1}]=0$:
$(S_k)$ is a martingale. By (a), $\E L_k^4\le c_4\ell_k^2$ with $\ell_k=\sum_{i<k}M_{ik}^2$; hence
$L_k$, $\gamma_k$ and $S_k$ have finite fourth moments. Expanding $(S_{k-1}+\gamma_k)^4$,
\[
\E S_k^4=\E S_{k-1}^4+4\,\E[S_{k-1}^3\gamma_k]+6\,\E[S_{k-1}^2\gamma_k^2]+4\,\E[S_{k-1}\gamma_k^3]+\E\gamma_k^4 .
\]
The second term is $8\,\E[S_{k-1}^3L_k]\,\E[u_k]=0$. For the fourth,
$4|S_{k-1}\gamma_k^3|\le2S_{k-1}^2\gamma_k^2+2\gamma_k^4$. So
$\E S_k^4\le\E S_{k-1}^4+8\,\E[S_{k-1}^2\gamma_k^2]+3\,\E\gamma_k^4$. By independence of $u_k$ from
$(S_{k-1},L_k)$ and the Cauchy--Schwarz inequality,
$\E[S_{k-1}^2\gamma_k^2]=4\,\E[S_{k-1}^2L_k^2]\le4(\E S_{k-1}^4)^{1/2}(\E L_k^4)^{1/2}$, and
$\E\gamma_k^4=16\,\kappa_k\,\E L_k^4$. Finally $\E S_{k-1}^4\le\E S_N^4=:\mathcal{M}$, by Jensen's inequality
applied to $S_{k-1}=\E[S_N\mid u_1,\dots,u_{k-1}]$. Summing over $k$, with $S_0=0$,
\[
\mathcal{M}\le32\,\mathcal{M}^{1/2}\sum_k(\E L_k^4)^{1/2}+48\,\bar\kappa\sum_k\E L_k^4
\le16\,c_4^{1/2}\|M\|^2\,\mathcal{M}^{1/2}+12\,\bar\kappa\,c_4\,\|M\|^4 ,
\]
because $\sum_k\ell_k=\frac12\|M\|^2$ and $\sum_k\ell_k^2\le(\sum_k\ell_k)^2$. An inequality
$x^2\le\alpha x+\beta$ with $x,\alpha,\beta\ge0$ implies $x\le\frac12(\alpha+\sqrt{\alpha^2+4\beta})$ and so
$x^2\le\alpha^2+2\beta$. With $x=\mathcal{M}^{1/2}$ this is the claim.

(e) Let $M^{o}$ be $M$ with its diagonal set to zero, so that $O=u'M^{o}u$, and use
$(D+O)^4\le8D^4+8O^4$. For the independent mean-zero summands $X_j=M_{jj}(u_j^2-1)$ of $D$,
$\E D^4=\sum_j\E X_j^4+3\sum_{j\ne\ell}\E X_j^2\,\E X_\ell^2\le\sum_j\E X_j^4+3(\sum_j\E X_j^2)^2$. Here
$\E X_j^2=M_{jj}^2(\kappa_j-1)\le\bar c^{1/2}M_{jj}^2$, because $\kappa_j\le(\E u_j^8)^{1/2}$; and
$(x-1)^4\le x^4+6x^2+1$ for $x\ge0$ gives $\E X_j^4\le M_{jj}^4(\bar c+6\bar c^{1/2}+1)\le8\bar c\,M_{jj}^4$,
because $\bar c\ge(\E u_j^2)^4=1$. So $\E D^4\le11\,\bar c\,(\sum_jM_{jj}^2)^2$. Apply (d) with
$\bar c^{1/2}$ as the bound on the fourth moments, so that $c_4$ becomes
$\max(3,\bar c^{1/2})\le3\bar c^{1/2}$: $\E O^4\le(256+24\bar c^{1/2})\,3\bar c^{1/2}\,\|M^{o}\|^4
\le840\,\bar c\,\|M^{o}\|^4$. Since
$(\sum_jM_{jj}^2)^2+\|M^{o}\|^4\le\|M\|^4$, the first claim follows with $8\cdot840=6720$. For
positive semidefinite $M$, $(u'Mu)^4\le8(u'Mu-\tr M)^4+8(\tr M)^4$ and $\|M\|\le\tr M$.
\end{proof}

\paragraph{Proof of the variance formula \eqref{eq:V}.} $\hat\theta-\mu'C\mu=2m'\varepsilon+\varepsilon'C\varepsilon=2a'u+u'Ku$. $K$ has zero diagonal, so by
Lemma~\ref{lem:forms}(b) the two terms are uncorrelated, with variances $4\|a\|^2$ and $2\|K\|^2$. Only second moments are used.\qed

\begin{lemma}[Moments of $\xi_U$]\label{lem:xi}
Under S1, S3 and S4, for every $U\in\mathcal{U}$: $1+\xi_U\ge0$;
\[
\E\xi_U=0,\qquad \E\xi_U^2=\nu_U,\qquad 0\le\nu_U\le\frac{\bnu}{d_U},\qquad
\E\xi_U^4\le\frac{c_Q}{d_U^2},\qquad \E(1+\xi_U)^4\le8\,(1+c_Q),
\]
with $c_Q=6720\,\bar c$ as in Lemma~\ref{lem:forms}(e).
\end{lemma}

\begin{proof}
By \eqref{eq:xi}, $1+\xi_U=u_U'\bar Q_Uu_U\ge0$ and $\xi_U=u_U'\bar Q_Uu_U-\tr\bar Q_U$, because
$\tr\bar Q_U=\tr\Pi_U/d_U=1$. Lemma~\ref{lem:forms}(b) with $\kappa_j=\kappa$ gives the mean and
$\E\xi_U^2=2\|\bar Q_U\|^2+(\kappa-3)\sum_j\bar Q_{U,jj}^{\,2}=\nu_U$. Since $\Pi_U$ is a projection,
$\|\bar Q_U\|^2=\tr\Pi_U/d_U^2=1/d_U$ and $0\le\Pi_{U,jj}\le1$, so
$\sum_j\bar Q_{U,jj}^{\,2}\le\sum_j\Pi_{U,jj}/d_U^2=1/d_U$; hence $\nu_U\le(2+|\kappa-3|)/d_U$. The
fourth moment is Lemma~\ref{lem:forms}(e) with $\|\bar Q_U\|^4=1/d_U^2$; this is where the eighth
moments of S4 are used. Finally $(1+\xi_U)^4\le8+8\xi_U^4$.
\end{proof}

\begin{lemma}[Mixed moments]\label{lem:mixed}
Under S1, S3 and S4, for every $U\in\mathcal{U}$,
\[
\E[\xi_U\,u_U]=0,\qquad \E[\xi_U\,u_Uu_U']=G_U,\qquad \|G_U\|_\op\le\frac{\bnu}{d_U}\le\bnu .
\]
Consequently, for non-random $H\in\mathbb{R}^{r\times|U|}$, $g\in\mathbb{R}^{r}$ and
$N\in\mathbb{R}^{|U|\times|U|}$,
\begin{equation}\label{eq:condid}
\E\big[\xi_U\,u_U'Nu_U\big]=\tr(NG_U),\qquad
\E\big[\xi_U\,\|g+Hu_U\|^2\big]=\tr\big(H'H\,G_U\big),
\end{equation}
and the second expression does not depend on $g$.
\end{lemma}

\begin{proof}
Write $\xi_U=\sum_{j,\ell}\bar Q_{U,j\ell}\,u_ju_\ell-1$, the sums running over $U$. For $p\in U$, in
$\E[\xi_Uu_p]=\sum_{j,\ell}\bar Q_{U,j\ell}\,\E[u_ju_\ell u_p]$ the single-index rule leaves only
$j=\ell=p$, and $\E u_p^3=0$ by S4. For $p\ne r$ in $U$, the single-index rule leaves in
$\E[\xi_Uu_pu_r]$ the two terms with $\{j,\ell\}=\{p,r\}$, which give $2\bar Q_{U,pr}$. For $p=r$ it
leaves the terms with $j=\ell$:
$\E[\xi_Uu_p^2]=\sum_{j\ne p}\bar Q_{U,jj}+\kappa\,\bar Q_{U,pp}-1=(\kappa-1)\bar Q_{U,pp}$, because
$\tr\bar Q_U=1$. Together this is $G_U=2\bar Q_U+(\kappa-3)\diag(\bar Q_U)$. The norm bound follows
from $\|\Pi_U\|_\op\le1$ and $0\le\Pi_{U,jj}\le1$. The first part of \eqref{eq:condid} is linearity;
for the second expand $\|g+Hu_U\|^2=\|g\|^2+2g'Hu_U+u_U'H'Hu_U$ and use $\E\xi_U=0$ and
$\E[\xi_Uu_U]=0$.
\end{proof}

Zero third moments enter through $\E[\xi_Uu_U]=0$ only. The second part of \eqref{eq:condid} is the
conditional-expectation identity that removes the own-set and mixed terms below.

\subsection{Propositions~\ref{prop:mean} and~\ref{prop:conc}}

We prove both for any $\kappa\ge1$ and any non-random $\Delta_i$ on the movers' rows. Step~0 is
Proposition~\ref{prop:mean}; it uses S1 to S3, the third and fourth moments of stayers' errors in S4, and second
moments of movers' errors, that is, A1 to A3. Steps 1 to 6 are Proposition~\ref{prop:conc} and use the moment
bounds of S4 and S5; the constants $c_1,c_2$ depend only on $\bar c$, $\bar\kappa$ and $\kappa$.

\begin{proof}[Proof of Propositions~\ref{prop:mean} and~\ref{prop:conc}]
All moments below are finite: by S4, S5 and Lemma~\ref{lem:forms} each piece of \eqref{eq:pieces} is
square integrable. Conditional expectations given a block of errors are computed by treating that
block as fixed, which S1 allows.

\medskip\noindent\emph{Step 0: means.} Since $w=a+Ku$ and $\E u=0$,
$\E P_1=4a'\Theta a+4\tr(K\Theta K)=4\sum_i\vartheta_i\big(a_i^2+\sum_\ell K_{i\ell}^2\big)$. For
$P_2$, fix $U$ and split
\begin{equation}\label{eq:gh}
w_U=g_U+h_U,\qquad h_U=K_{UU}\,u_U,\qquad g_U=a_U+K_{U,-U}\,u_{-U},
\end{equation}
where $-U$ denotes all rows not in $U$. Given $u_{-U}$, $g_U$ is fixed and $u_U$ keeps its
distribution, so by \eqref{eq:condid}
\begin{equation}\label{eq:eU}
\E\big[\xi_U\|w_U\|^2\,\big|\,u_{-U}\big]=\tr\big(K_{UU}^2G_U\big)=:e_U ,
\end{equation}
a constant; thus $\E P_2=4\sum_Ue_U$. For $P_3$, $\E[(1+\xi_i)(1+\xi_\ell)]$ equals $1+\nu_U$ if $i$
and $\ell$ lie in the same set $U$ and $1$ otherwise (Lemma~\ref{lem:xi} and independence across
sets), so $\E P_3=-2\sum_{i,\ell}K_{i\ell}^2\vartheta_i\vartheta_\ell-2\sum_U\nu_Uk_{UU}$. Finally
$\E P_4=-\sum_U\tau_U^2\chi_U$. By Lemma~\ref{lem:stayer}, $4e_U-2\nu_Uk_{UU}-\tau_U^2\chi_U=0$ for every $U$, so
\[
\E\hat V=4\sum_i\vartheta_ia_i^2+4\sum_{i,\ell}\vartheta_iK_{i\ell}^2-2\sum_{i,\ell}K_{i\ell}^2\vartheta_i\vartheta_\ell .
\]
Put $\vartheta_i=1+\delta_i$ with $\delta_i=\Delta_i/\sigma_i^2$ on $R$ and $\delta_i=0$ elsewhere. The
terms linear in $\delta$ from the two double sums cancel, which leaves
$4\|a\|^2+2\|K\|^2+4\sum_{i\in R}\delta_ia_i^2-2\sum_{i,\ell\in R}K_{i\ell}^2\delta_i\delta_\ell$. By
\eqref{eq:V}, and since $\delta_ia_i^2=\Delta_im_i^2$ and
$K_{i\ell}^2\delta_i\delta_\ell=C_{i\ell}^2\Delta_i\Delta_\ell$, this is $V+b(\Delta)$, which is Proposition~\ref{prop:mean}.

\medskip\noindent For Proposition~\ref{prop:conc} we bound the variance of each piece and use
\begin{equation}\label{eq:four}
\Var(\hat V)\le4\sum_{k=1}^4\Var(P_k),
\end{equation}
which is $(\sum_{k=1}^4x_k)^2\le4\sum_kx_k^2$ applied to the centred pieces. This disposes of all
covariances between different pieces.

\medskip\noindent\emph{Step 1: piece (1).} $P_1-\E P_1=8\,(K\Theta a)'u+4\,\big(u'K\Theta Ku-\tr(K\Theta K)\big)$.
By Lemma~\ref{lem:forms}(a),(b) and $\Var(X+Y)\le2\Var X+2\Var Y$,
\[
\Var(P_1)\le128\,\|K\Theta a\|^2+32\max(2,\bar\kappa-1)\,\|K\Theta K\|^2 .
\]
This uses fourth moments of all rows and no third moment: the covariance between the linear
and the quadratic part, which involves third moments, is not computed but bounded.

\medskip\noindent\emph{Step 2: piece (2), terms of the same set.} Let $Z_U=\xi_U\|w_U\|^2-e_U$, so that
$P_2-\E P_2=4\sum_UZ_U$ and
\begin{equation}\label{eq:S2}
\Var(P_2)=16\sum_U\E Z_U^2+16\sum_{U\ne U'}\E[Z_UZ_{U'}] .
\end{equation}
With \eqref{eq:gh}, $\|w_U\|^2=\|g_U\|^2+2g_U'h_U+\|h_U\|^2\le2\|g_U\|^2+2\|h_U\|^2$, so
\[
\E Z_U^2\le\E\big[\xi_U^2\|w_U\|^4\big]\le8\,\E\big[\xi_U^2\|g_U\|^4\big]+8\,\E\big[\xi_U^2\|h_U\|^4\big].
\]
\emph{First term.} $g_U$ is independent of $\xi_U$, so it equals $8\,\nu_U\,\E\|g_U\|^4$. By
Lemma~\ref{lem:forms}(c), applied to the errors outside $U$ (fourth moments, S4 and S5), and
$\nu_U\le\bnu$, it is at most $32\,\bnu\,c_4\,(\|a_U\|^2+\|K_{U,-U}\|^2)^2$.
\emph{Second term.} $\|h_U\|^2=u_U'K_{UU}^2u_U$ with $K_{UU}^2$ positive semidefinite and
$\tr K_{UU}^2=k_{UU}$. By the Cauchy--Schwarz inequality, Lemma~\ref{lem:xi} and
Lemma~\ref{lem:forms}(e) (eighth moments of the rows of $U$),
\[
\E\big[\xi_U^2\|h_U\|^4\big]\le\big(\E\xi_U^4\big)^{1/2}\big(\E(u_U'K_{UU}^2u_U)^4\big)^{1/2}
\le c_Q^{1/2}\,\big(8(c_Q+1)\big)^{1/2}k_{UU}^2\le3\,(c_Q+1)\,k_{UU}^2 .
\]
Since $\|a_U\|^2+\|K_{U,-U}\|^2+k_{UU}=\rho_U$,
\begin{equation}\label{eq:within}
\E Z_U^2\le c_{2a}\,\rho_U^2,\qquad c_{2a}=32\,\bnu\,c_4+24\,(c_Q+1) .
\end{equation}

\medskip\noindent\emph{Step 3: piece (2), terms of two different sets.} Fix $U\ne U'$, write
$\Lambda=K_{UU'}$ (so $K_{U'U}=\Lambda'$) and let $O$ be the set of rows outside $U\cup U'$. The three blocks
$u_U$, $u_{U'}$, $u_O$ are independent. Put $\bar g_U=a_U+K_{UO}u_O$ and
$\bar g_{U'}=a_{U'}+K_{U'O}u_O$, so that
$w_U=\bar g_U+h_U+\Lambda u_{U'}$ and $w_{U'}=\bar g_{U'}+h_{U'}+\Lambda'u_U$. Split $Z_U=Z_U^0+Z_U^1$ with
\[
Z_U^0=\xi_U\|\bar g_U+h_U\|^2-e_U,\qquad
Z_U^1=\underbrace{2\,\xi_U\,\bar g_U'\Lambda u_{U'}}_{Z_U^{1a}}+\underbrace{2\,\xi_U\,u_U'K_{UU}\Lambda u_{U'}}_{Z_U^{1b}}
+\underbrace{\xi_U\,u_{U'}'\Lambda'\Lambda u_{U'}}_{Z_U^{1c}},
\]
and likewise $Z_{U'}=Z_{U'}^0+Z_{U'}^1$ with the roles of $U$ and $U'$ exchanged and $\Lambda'$ in place of
$\Lambda$: $Z_{U'}^{1a}=2\xi_{U'}\bar g_{U'}'\Lambda'u_U$, $Z_{U'}^{1b}=2\xi_{U'}u_{U'}'K_{U'U'}\Lambda'u_U$ and
$Z_{U'}^{1c}=\xi_{U'}u_U'\Lambda\Lambda'u_U$.
$Z_U^0$ is a function of $(u_U,u_O)$ and $Z_{U'}^0$ of $(u_{U'},u_O)$.

\emph{(3a) $\E[Z_U^0Z_{U'}^0]=0$.} Given $(u_O,u_{U'})$, $Z_{U'}^0$ and $\bar g_U$ are fixed, and by
\eqref{eq:condid} $\E[\xi_U\|\bar g_U+h_U\|^2\mid u_O,u_{U'}]=e_U$. So
$\E[Z_U^0\mid u_O,u_{U'}]=0$.

\emph{(3b) $\E[Z_U^0Z_{U'}^1]=0$ and $\E[Z_U^1Z_{U'}^0]=0$.} Given $(u_O,u_U)$, $Z_U^0$, $\bar g_{U'}$
and $\Lambda'u_U$ are fixed, and
$\E[Z_{U'}^1\mid u_O,u_U]=2\,\E[\xi_{U'}]\,\bar g_{U'}'\Lambda'u_U+2\,\E[\xi_{U'}u_{U'}]'K_{U'U'}\Lambda'u_U
+\E[\xi_{U'}]\,\|\Lambda'u_U\|^2=0$ by Lemmas \ref{lem:xi} and \ref{lem:mixed}. The other statement is symmetric.

\emph{(3c) $\E[Z_U^1Z_{U'}^1]$.} Each of the nine products is, given $u_O$, a finite sum of terms
(factor in $u_U$)$\,\times\,$(factor in $u_{U'}$), and the two blocks are independent. The factor in
$u_U$ is $\xi_U$ times a polynomial in $u_U$, and it has mean zero whenever that polynomial is linear
($\E[\xi_Uu_U]=0$); the same holds for $U'$. The table lists, for each product, the degree in $u_U$ of the polynomial that multiplies $\xi_U$,
the degree in $u_{U'}$ of the one that multiplies $\xi_{U'}$, and the expectation.
\begin{center}\small
\renewcommand{\arraystretch}{1.25}
\begin{tabular}{@{}l|ccc@{}}
 & $Z_{U'}^{1a}$ & $Z_{U'}^{1b}$ & $Z_{U'}^{1c}$ \\
\hline
$Z_U^{1a}$ & $(1,1)$: $0$ & $(1,2)$: $0$ & $(2,1)$: $0$ \\
$Z_U^{1b}$ & $(2,1)$: $0$ & $(2,2)$: $4\tr\big(K_{UU}\Lambda G_{U'}K_{U'U'}\Lambda'G_U\big)$ & $(3,1)$: $0$ \\
$Z_U^{1c}$ & $(1,2)$: $0$ & $(1,3)$: $0$ & $(2,2)$: $\tr(\Lambda'\Lambda G_{U'})\,\tr(\Lambda\Lambda'G_U)$ \\
\end{tabular}
\end{center}
For the entry $(b,b)$ take the expectation over $u_{U'}$ first. Because
$\E[\xi_{U'}u_{U'}u_{U'}']=G_{U'}$, this leaves the quadratic form
$4\,\xi_U\,u_U'\big(K_{UU}\Lambda G_{U'}K_{U'U'}\Lambda'\big)u_U$, and \eqref{eq:condid} gives the trace. The
entry $(c,c)$ is \eqref{eq:condid} applied twice. Every
product is integrable: it is a polynomial of degree at most five in the errors of each of the two
sets (S4 gives eight moments) and of degree at most two in $u_O$. No fifth or third moment is
evaluated, because the factors of degree three are multiplied by a factor with mean zero.
Altogether, exactly,
\begin{equation}\label{eq:cross}
\E[Z_UZ_{U'}]=4\tr\big(K_{UU}\,\Lambda\,G_{U'}\,K_{U'U'}\,\Lambda'\,G_U\big)+\tr\big(\Lambda'\Lambda\,G_{U'}\big)\tr\big(\Lambda\Lambda'\,G_U\big),
\qquad \Lambda=K_{UU'} .
\end{equation}
By $|\tr(XY)|\le\|X\|\|Y\|$, $\|XYZ\|\le\|X\|\|Y\|\|Z\|_\op$ and $\|G_U\|_\op\le\bnu$
(Lemma~\ref{lem:mixed}), the first term is at most
$4\bnu^2\,k_{UU}^{1/2}k_{U'U'}^{1/2}k_{UU'}\le2\bnu^2(k_{UU}+k_{U'U'})k_{UU'}$ in absolute value; since
$\Lambda'\Lambda$ and $\Lambda\Lambda'$ are positive semidefinite with trace $k_{UU'}$, the second is at most
$\bnu^2k_{UU'}^2$. Summing over ordered pairs with
\eqref{eq:facts},
\[
\sum_{U\ne U'}\big|\E[Z_UZ_{U'}]\big|\le\bnu^2\Big(4\sum_Uk_{UU}\sum_{U'}k_{UU'}+\sum_U\Big(\sum_{U'}k_{UU'}\Big)^2\Big)
\le5\,\bnu^2\sum_U\rho_U^2 .
\]
With \eqref{eq:S2} and \eqref{eq:within}, $\Var(P_2)\le16\,(c_{2a}+5\bnu^2)\sum_U\rho_U^2$.

\medskip\noindent\emph{Step 4: piece (3).} Let
$\tilde k_U=\sum_{i\in U}\sum_{\ell\in R}K_{i\ell}^2\vartheta_\ell+\sum_{U'\in\mathcal{U}}k_{UU'}$, so that
$|\tilde k_U|\le\bth\,k_U$. Expanding $(1+\xi_i)(1+\xi_\ell)$ in \eqref{eq:pieces} and collecting,
\[
P_3-\E P_3=X+Y,\qquad X=\sum_U\Big[-4\,\tilde k_U\,\xi_U-2\,k_{UU}\,(\xi_U^2-\nu_U)\Big],\qquad
Y=-2\sum_{U\ne U'}k_{UU'}\,\xi_U\xi_{U'} .
\]
$X$ is the linear part together with the same-set squares, and $Y$ the cross-set products. The summands of $X$ are independent and centred, so by Lemma~\ref{lem:xi}
$\E X^2\le\sum_U\big(32\,\tilde k_U^2\,\nu_U+8\,k_{UU}^2\,\E\xi_U^4\big)\le(32\,\bnu\,\bth^2+8\,c_Q)\sum_Uk_U^2$;
the same-set squares are where $\E\xi_U^4$, hence eighth moments, are needed. The products
$\xi_U\xi_{U'}$ over unordered pairs $U\ne U'$ have mean zero and are uncorrelated with one another
(two different pairs differ in a set, whose $\xi$ is independent of the other factors and has mean zero), so
$\E Y^2=8\sum_{U\ne U'}k_{UU'}^2\nu_U\nu_{U'}\le8\,\bnu^2\sum_Uk_U^2$ by \eqref{eq:facts}. Hence
$\Var(P_3)\le2\E X^2+2\E Y^2\le(64\,\bnu\,\bth^2+16\,c_Q+16\,\bnu^2)\sum_U\rho_U^2$.

\medskip\noindent\emph{Step 5: piece (4).} The summands are independent, $1+\nu_U\ge1$, and by
Lemmas \ref{lem:stayer} and \ref{lem:xi}, $\tau_U^2|\chi_U|\le\bnu\,k_{UU}$ and $\E(1+\xi_U)^4\le8(1+c_Q)$. So
\[
\Var(P_4)=\sum_U\frac{\tau_U^4\chi_U^2}{(1+\nu_U)^2}\Var\big((1+\xi_U)^2\big)\le128\,\bnu^2\,(1+c_Q)\sum_U\rho_U^2 .
\]

\medskip\noindent\emph{Step 6: assembly.} By \eqref{eq:four} and Steps 1 to 5 (recall $\bth\ge1$),
$\Var(\hat V)\le c_1\big(\|K\Theta a\|^2+\|K\Theta K\|^2\big)+c_2\,\bth^2\sum_U\rho_U^2$ with
\[
c_1=128\max(4,\bar\kappa-1),\qquad
c_2=64\,(c_{2a}+5\bnu^2)+4\,(64\,\bnu+16\,c_Q+16\,\bnu^2)+512\,\bnu^2\,(1+c_Q).
\]
Next, $\|K\Theta a\|\le\lambda\,\bth\,\|a\|$ and
$\|K\Theta K\|\le\lambda\,\bth\,\|K\|$, so the first bracket is at most
$\lambda^2\bth^2(\|a\|^2+\|K\|^2)\le\frac12\lambda^2\bth^2V$; and by \eqref{eq:facts}
$\sum_U\rho_U^2\le\max_U\rho_U\sum_U\rho_U\le\frac14\max_U\phi_U\,V^2$. Hence
$\Var(\hat V)\le\bth^2\big(\frac12c_1\lambda^2V+\frac14c_2\max_U\phi_UV^2\big)$; dividing by $V^2$ gives
Proposition~\ref{prop:conc} with $c'=\max(\frac12c_1,\frac14c_2)$, because $\E\hat V=V+b(\Delta)$.
\end{proof}

\paragraph{Departures from A2 and A3.} Three calculations give the mean of the estimator \eqref{eq:Vhat} of the
text (the adjustment for $\kappa=3$) when stayers' errors depart from A2 or A3 and movers' rows have their true
variances.

\emph{Another kurtosis.} If stayers' errors have kurtosis $\kappa\ne3$, Step~0 gives
$\E\hat V-V=\sum_U\tau_U^2\big[\chi_U(\kappa)-\chi_U(3)\,(1+\nu_U(\kappa))/(1+\nu_U(3))\big]$, where the arguments
indicate the kurtosis used in \eqref{eq:nuh}. Inserting \eqref{eq:nuh},
\begin{equation}\label{eq:kurt}
\E\hat V-V=2\,(\kappa-3)\sum_U\tau_U^2\,c_U^2\,\frac{(|U|-1)^2(|U|+4)}{|U|\,(|U|+1)} .
\end{equation}

\emph{Unequal variances within a stayer.} Let the rows of $U$ have normal errors with variances $v_1,\dots,v_t$,
$t=|U|$, with mean $\bar v$ and $s_U^2=\frac{1}{t-1}\sum_j(v_j-\bar v)^2$. The rows of $U$ have a common $m_i$ and
common weights to every row outside $U$, so every term of $V$ and of $\E\hat V$ that involves one row of $U$ at a
time depends on the $v_j$ through $\sum_jv_j$ only, and $\E q_U=\bar v$: these terms agree. The terms within $U$ do
not. With $S=\sum_j\varepsilon_j$ and $Q=\sum_j\varepsilon_j^2$ over $U$ we have $q_U=(Q-S^2/t)/(t-1)$ and
$\sum_{i\in U}(S-\varepsilon_i)^2=(t-2)S^2+Q$. The normal moments $\E S^4=3(\sum_jv_j)^2$ and
$\E Q^2=\E[S^2Q]=(\sum_jv_j)^2+2\sum_jv_j^2$ then give, after collecting terms,
\begin{equation}\label{eq:uneq}
\E\hat V-V=6\sum_Uc_U^2\,s_U^2\,\frac{(|U|-1)^2(|U|+4)}{|U|\,(|U|+1)}\;\ge0 .
\end{equation}

\emph{Skewness.} If S4's $\E\varepsilon_i^3=0$ is replaced by $\E\varepsilon_i^3=\mu_{3U}$ on $U$,
only the mean of $P_2$ changes: $\E[\xi_Uu_p]=\E u_p^3/|U|$ for $p\in U$, so
$\E[\xi_U\,2g_U'h_U]=2\,\E[g_U]'K_{UU}\E[\xi_Uu_U]=2\,a_U'K_{UU}\mathbf{1}\,\E u^3/|U|$ in the notation of
\eqref{eq:gh}. With $a_U=\tau_U^{1/2}m_U\mathbf{1}$, $K_{UU}\mathbf{1}=\tau_Uc_U(|U|-1)\mathbf{1}$ and
$\E u^3=\mu_{3U}/\tau_U^{3/2}$ this adds $8\,(|U|-1)\,c_U\,m_U\,\mu_{3U}$ to $\E\hat V$ for each stayer, as stated
after Proposition~\ref{prop:mean}.

\subsection{Proposition~\ref{prop:fitted} and Corollary~\ref{cor:transfer}}

Now the regression structure matters. $X$ is the $n\times k$ matrix with rows $x_i'$, of full column rank,
$P=XS_{xx}^{-1}X'$, $M=I-P$, and $\hat\sigma_i^2=y_i\,(My)_i/M_{ii}$. Assumption~1 of KSS is used in this form:
\begin{enumerate}[label=\textbf{(K\arabic*)},leftmargin=3.3em,itemsep=3pt]
\item for all rows, $\E\varepsilon_i^4\le\bar\mu_4$, $\sigma_i^2\ge\underline\sigma^2>0$, $P_{ii}\le\bar P<1$ and
$|x_i'\beta|\le\bar\mu$. Consequently $\sigma_i^2\le\bar\sigma^2:=\bar\mu_4^{1/2}$, $M_{ii}\ge1-\bar P$, and S5's
moment bound holds with $\bar\kappa=\bar\mu_4/\underline\sigma^4$.
\end{enumerate}

\begin{lemma}[The KSS variance estimates]\label{lem:kssvar}
Under S1 and K1, $\E\hat\sigma_i^2=\sigma_i^2$ for every $i$, and for every $t\in\mathbb{R}^n$
\[
\Var\Big(\sum_it_i\,\hat\sigma_i^2\Big)\le c_\sigma\,\|t\|^2,\qquad
c_\sigma=\frac{2\,\bar\sigma^2\bar\mu^2}{(1-\bar P)^2}+\frac{4\,\bar\sigma^4}{1-\bar P}+2\,\bar\mu_4 .
\]
That is, the covariance matrix of $(\hat\sigma_1^2,\dots,\hat\sigma_n^2)$ has operator norm at most $c_\sigma$.
\end{lemma}

\begin{proof}
$MX=0$ gives $My=M\varepsilon$, so
$\hat\sigma_i^2=\mu_i(M\varepsilon)_i/M_{ii}+\varepsilon_i(M\varepsilon)_i/M_{ii}$. The first part has
mean zero and the second has mean $M_{ii}\sigma_i^2/M_{ii}$. Next,
$\sum_it_i\hat\sigma_i^2=\ell'\varepsilon+\varepsilon'\Gamma\varepsilon$ with $\ell=Mr$,
$r_i=t_i\mu_i/M_{ii}$, and $\Gamma=\frac12(D_tM+MD_t)$, $D_t=\diag(t_i/M_{ii})$.
\emph{Linear part.} $\Var(\ell'\varepsilon)=\sum_j\ell_j^2\sigma_j^2\le\bar\sigma^2\|Mr\|^2\le\bar\sigma^2\|r\|^2
\le\bar\sigma^2\bar\mu^2(1-\bar P)^{-2}\|t\|^2$, because $M$ is a projection.
\emph{Quadratic part.} Apply Lemma~\ref{lem:forms}(b) to $u_j=\varepsilon_j/\sigma_j$ and the matrix
with entries $\sigma_j\Gamma_{j\ell}\sigma_\ell$:
\[
\Var(\varepsilon'\Gamma\varepsilon)=2\sum_{j,\ell}\Gamma_{j\ell}^2\sigma_j^2\sigma_\ell^2+\sum_j\big(\E\varepsilon_j^4-3\sigma_j^4\big)\Gamma_{jj}^2
\le2\,\bar\sigma^4\,\|\Gamma\|^2+\bar\mu_4\sum_j\Gamma_{jj}^2 .
\]
Here $\Gamma_{jj}=t_j$, and $\|\Gamma\|\le\|D_tM\|$ with
$\|D_tM\|^2=\sum_i(t_i/M_{ii})^2\sum_\ell M_{i\ell}^2=\sum_it_i^2/M_{ii}\le\|t\|^2/(1-\bar P)$, because
$\sum_\ell M_{i\ell}^2=(M^2)_{ii}=M_{ii}$. So
$\Var(\varepsilon'\Gamma\varepsilon)\le(2\bar\sigma^4/(1-\bar P)+\bar\mu_4)\|t\|^2$; the uniform bound
on fourth moments in KSS Assumption~1(i) is used here. The claim follows from
$\Var(X+Y)\le2\Var X+2\Var Y$.
\end{proof}

\begin{proof}[Proof of Proposition~\ref{prop:fitted}]
Write $\hat V(v)$ for \eqref{eq:VhatA} when the movers' rows receive the values $v=(v_i)_{i\in R}$, and let
$v^*_i=\psi_i(\eta^*)$, $\hat v_i=\st_i^2$ the fitted values, $e_i=\hat v_i-v^*_i$ and $\delta=\hat\eta-\eta^*$. Let
$\Omega$ be the event that $\|\delta\|\le r$ and $\hat v_i=\psi_i(\hat\eta)$ for all $i\in R$. By assumption and
Markov's inequality, $\Pr(\Omega^c)\to0$ as $n_0\to\infty$.

\emph{Step A: an identity.} $\hat V(v)$ is a polynomial of degree two in $v$: the first term of
\eqref{eq:VhatA} is linear in the movers' values, the second is quadratic, and the third does not involve them.
Expanding around $v^*$, with $C$ symmetric,
\begin{equation}\label{eq:expand}
\hat V(\hat v)-\hat V(v^*)=4\sum_{i\in R}e_i\Big[\hat m_i^{\,2}-\sum_\ell C_{i\ell}^2\,\st_\ell^{*2}\Big]
-2\sum_{i\in R}\sum_{\ell\in R}C_{i\ell}^2\,e_ie_\ell ,
\end{equation}
where $\st_\ell^{*2}=q_U$ for $\ell\in U$ and $\st_\ell^{*2}=v^*_\ell$ for $\ell\in R$. This is an algebraic
identity, valid for every realisation, whatever the dependence between $\hat\eta$ and the other terms.

\emph{Step B: the brackets.} Let $J=\sum_{i\in R}\hat m_i^{\,2}+\sum_{i\in R}\sum_\ell C_{i\ell}^2|\st_\ell^{*2}|$.
Since $C_{i\ell}^2=K_{i\ell}^2/(\sigma_i^2\sigma_\ell^2)$ and $\sigma_i^2\ge\underline\sigma^2$, we have
$\sum_{i,\ell\in R}C_{i\ell}^2\le\underline\sigma^{-4}\|K\|^2\le\frac12\underline\sigma^{-4}V$, and
$J\le\underline\sigma^{-2}\big[\|w_R\|^2+\sum_{i\in R}\sum_\ell K_{i\ell}^2\theta^*_\ell\big]$ with
$\theta^*_\ell=|\st_\ell^{*2}|/\sigma_\ell^2$, which is $1+\xi_U$ on a set and at most
$\bth:=\max(1,\bar\psi/\underline\sigma^2)$ on $R$. By Lemma~\ref{lem:forms}(c),
$(\E\|w_R\|^4)^{1/2}\le2c_4^{1/2}(\|a_R\|^2+\|K_{R\cdot}\|^2)$; by Minkowski's inequality and
$\E(1+\xi_U)^2=1+\nu_U$, the second sum has $L^2$ norm at most $\max(\bth,(1+\bnu)^{1/2})\,\|K_{R\cdot}\|^2$. As
$\|a_R\|^2+\|K_{R\cdot}\|^2\le\frac12V$,
\[
\big(\E J^2\big)^{1/2}\le c_J\,V,\qquad c_J=\frac{1}{2\,\underline\sigma^2}\Big(2c_4^{1/2}+\bth+(1+\bnu)^{1/2}\Big).
\]
\emph{Step C: the fitted values.} On the event $\Omega$ we have $|e_i|\le L\|\delta\|$ for all
$i\in R$, so by \eqref{eq:expand}
$|\hat V(\hat v)-\hat V(v^*)|\le4L\|\delta\|\,J+2L^2\|\delta\|^2\sum_{i,\ell\in R}C_{i\ell}^2$. By the
Cauchy--Schwarz inequality and $\E\|\delta\|^2\le c_\eta/n_0$,
\[
\E\Big[\big|\hat V(\hat v)-\hat V(v^*)\big|\,;\,\Omega\Big]\le V\Big(4L\,c_J\,c_\eta^{1/2}\,n_0^{-1/2}+L^2\underline\sigma^{-4}\,c_\eta\,n_0^{-1}\Big),
\]
so that $[\hat V(\hat v)-\hat V(v^*)]/V=O_p(n_0^{-1/2})$, because $\Pr(\Omega^c)\to0$.

\emph{Step D: the estimator at $v^*$.} $\hat V(v^*)$ is the estimator of Propositions~\ref{prop:mean}
and~\ref{prop:conc} with the non-random $\Delta_i=v^*_i-\sigma_i^2$, for which $|\vartheta_i|\le\bth$. Hence
$\E\hat V(v^*)=V+b(\Delta)$ and
$\E|\hat V(v^*)-V-b(\Delta)|\le c'^{1/2}\,\bth\,(\max_U\phi_U+\lambda_{\max}^2/V)^{1/2}\,V$. Steps C and D and
Markov's inequality give the proposition.
\end{proof}

\begin{proof}[Proof of Corollary~\ref{cor:transfer}]
Let $\|b_i\|\le\bar b$ and $\|z_i\|\le\bar z$, let $\underline\omega>0$ bound the two smallest eigenvalues from
below, and write $G_z=n_{\cR}^{-1}\sum_{i\in\cR}z_iz_i'$, which is non-random. The statistic is
$\hat\eta=(\hat N_1,\hat N_2)$ with
\[
\hat N_1=\frac{1}{n_{\cR}}\sum_{i\in\cR}z_i\,\hat\sigma_i^2,\qquad
\hat N_2=\frac{1}{n_{\cR}}\sum_{i\in\cR}z_ib_i'\,\hat\sigma_i^2,
\]
a vector and a matrix of fixed dimensions, and $\eta^*=(N_1^*,N_2^*)$ is the same with $\sigma_i^2$ in place of
$\hat\sigma_i^2$.

\emph{Rate.} Lemma~\ref{lem:kssvar}, applied to each entry of $\hat N_1$ and $\hat N_2$ (with $t_i=z_{ij}/n_{\cR}$ or
$t_i=z_{ij}b_{ik}/n_{\cR}$ on $\cR$ and zero elsewhere) and summed, gives
$\E\|\hat N_1-N_1^*\|^2\le c_\sigma\bar z^2/n_{\cR}$ and $\E\|\hat N_2-N_2^*\|^2\le c_\sigma\bar z^2\bar b^2/n_{\cR}$. So
$\E\|\hat\eta-\eta^*\|^2\le c_\eta/n_0$ with $n_0=n_{\cR}$ and $c_\eta=c_\sigma\bar z^2(1+\bar b^2)$. The leave-out
estimates of different rows need not be independent for this.

\emph{The fitted values as a function of $\eta$.} Without floors, the two steps are the following maps of
$\eta=(N_1,N_2)$. Step~1: $c_z=G_z^{-1}N_1$ and $s_i=z_i'c_z$. Step~2: with $x_i=s_i\,b_i$, the coefficient is
$H(\eta)^{-1}\varrho(\eta)$, where
\[
H(\eta)=\frac{1}{n_{\cR}}\sum_{i\in\cR}x_ix_i',\qquad \varrho(\eta)=N_2'c_z ,
\]
because $n_{\cR}^{-1}\sum_is_ib_i\hat\sigma_i^2=\hat N_2'c_z$ at $N_2=\hat N_2$. The fitted value is
\[
\psi^\circ_i(\eta)=x_i(\eta)'\,H(\eta)^{-1}\varrho(\eta) .
\]
At $\eta^*$ it equals $v_i^*$: $H(\eta^*)^{-1}\varrho(\eta^*)$ is the coefficient $d^*$ of the regression of
$\sigma_i^2$ on $s_i^*b_i$. The maps $\eta\mapsto s_i,x_i,H,\varrho$ are polynomials in $(N_1,N_2)$ with bounded
coefficients, and the inverse is handled by $H^{-1}-H_*^{-1}=H^{-1}(H_*-H)H_*^{-1}$;
$\|G_z^{-1}\|_\op\le1/\underline\omega$, the smallest eigenvalue of $H(\eta^*)$ is at least $\underline\omega$, and
$\|N_1^*\|$, $\|N_2^*\|$ are bounded by $\bar z\bar\sigma^2$ and $\bar z\bar b\bar\sigma^2$. So there are $r_1>0$ and
$L$, depending only on these constants, such that for $\|\eta-\eta^*\|\le r_1$ the smallest eigenvalue of
$H(\eta)$ is at least $\underline\omega/2$ and the maps $\eta\mapsto s_i$ and $\eta\mapsto\psi^\circ_i(\eta)$ are
Lipschitz with constant $L$ for every $i$.

\emph{Floors.} The estimator floors $\hat s_i$ at $\hat f_s=0.05\,n_{\cR}^{-1}\sum_{i\in\cR}\hat\sigma_i^2$, which is
$0.05$ times the first entry of $\hat N_1$, and the fitted value at a random number $\hat f_v\ge0$. When a floor
binds for some row the fitted values are no longer the function $\psi^\circ$ of $\hat\eta$; this is why
Proposition~\ref{prop:fitted} asks for the representation $\st_i^2=\psi_i(\hat\eta)$ only with probability tending
to one. By assumption there is a margin $\iota>0$ with $s_i^*\ge f_s^*+\iota$ for all $i\in\cR$, where $f_s^*$ is
$0.05$ times the first entry of $N_1^*$, and $\Pr\big(\hat f_v\le\min_{i\in\cR}v_i^*-\iota\big)\to1$. Put
$\psi_i=\psi_i^\circ$ on $\{\|\eta-\eta^*\|\le r_1\}$ and $\psi_i=0$ elsewhere, and
$r=\min\big(r_1,\iota/(2L+2)\big)$. If $\|\hat\eta-\eta^*\|\le r$, then $\hat s_i-\hat f_s\ge\iota-(L+1)r>0$, so
the first floor binds for no row, and the fitted value before the last floor is
$\psi_i^\circ(\hat\eta)\ge v_i^*-\iota/2$; on the event in the last display it exceeds $\hat f_v$, so the last
floor does not bind either. Hence $\st_i^2=\psi_i(\hat\eta)$ for all $i\in\cR$ with probability tending to one,
and the conditions of Proposition~\ref{prop:fitted} hold, with $\bar\psi$ a bound on $|v_i^*|$ that follows from
the bounds above. Floors that keep binding in large samples are not covered. If $\sigma_i^2=s_i^*\,b_i'd$ on
$\cR$, the regression of $\sigma_i^2$ on $s_i^*b_i$ fits exactly, so $v_i^*=\sigma_i^2$ and $\Delta=0$. If
$\sigma_i^2=(z_i'c)(b_i'd)$ and $\sum_{i\in\cR}z_i\,(z_i'c)\,(b_i'd-\bar\pi)=0$ for a constant $\bar\pi\neq0$, then
$c_z^*=\bar\pi\,c$, so $s_i^*$ is proportional to $z_i'c$ and the first case applies.
\end{proof}

\subsection{Proposition~\ref{prop:sim}: the simulated double sum}

Let $\eta^{(1)},\dots,\eta^{(D_s)}$ be independent $N(0,I_n)$ vectors, drawn independently of the data, and
$\zeta^{(d)}_\ell=\st_\ell\,\eta^{(d)}_\ell$. $\hat V_{D_s}$ is \eqref{eq:VhatA} with the double sum replaced by
$\frac1{D_s}\sum_d\sum_i\st_i^2\big(\sum_\ell C_{i\ell}\zeta^{(d)}_\ell\big)^2$. We show, under S1 to S6 and
$0\le\vartheta_i\le\bth$,
\[
\E\big[(\hat V_{D_s}-\hat V)^2\big]=\frac{8}{D_s}\,\E\tr\big(\tilde K^4\big)\le\frac{c_3\,\bth^4}{D_s}\,\lambda_{\max}^2\,\|K\|^2,
\qquad \tilde K_{i\ell}=\st_iC_{i\ell}\st_\ell,\quad c_3=16(1+\bnu)+48(1+c_Q)+32\,\bnu^2 ,
\]
which gives the proposition with $c''=c_3/2$ because $\|K\|^2\le\frac12V$. In the proof $D$ stands for $D_s$.

\begin{proof}
By construction,
$\sum_i\st_i^2(\sum_\ell C_{i\ell}\zeta^{(d)}_\ell)^2=\|\tilde K\eta^{(d)}\|^2=\eta^{(d)\prime}\tilde K^2\eta^{(d)}$,
which given the data has mean $\tr\tilde K^2=\sum_{i,\ell}C_{i\ell}^2\st_i^2\st_\ell^2$ and variance
$2\tr\tilde K^4$ (Lemma~\ref{lem:forms}(b) with $\kappa_j=3$). The double sum enters with the factor
$-2$, so $\hat V_D-\hat V$ has conditional mean zero and conditional variance $8\tr(\tilde K^4)/D$.
This gives the equality.

For the bound let $\theta_i=\st_i^2/\sigma_i^2=\vartheta_i(1+\xi_i)\ge0$, so that
$\tilde K=\diag(\theta_i)^{1/2}K\diag(\theta_i)^{1/2}$ and
$\tr\tilde K^4=\|\tilde K^2\|^2=\sum_{i,k}\theta_i\theta_k\,(K\diag(\theta)K)_{ik}^2$. Since $\vartheta_j=1$ on
sets and $\xi_j=0$ on $R$, $\diag(\theta)=\Theta+\diag(\xi)$ and
$K\diag(\theta)K=K\Theta K+\sum_U\xi_UW^U$ with $W^U=K_{\cdot U}K_{U\cdot}$, a symmetric $n\times n$
matrix. By $(x+y)^2\le2x^2+2y^2$,
\begin{equation}\label{eq:tr4}
\tr\tilde K^4\le2\sum_{i,k}\theta_i\theta_k\,(K\Theta K)_{ik}^2+2\sum_{i,k}\theta_i\theta_k\Big(\sum_U\xi_UW^U_{ik}\Big)^2 .
\end{equation}
\emph{First sum.} $\E[\theta_i\theta_k]=\vartheta_i\vartheta_k\,\E[(1+\xi_i)(1+\xi_k)]\le\bth^2(1+\bnu)$, so its
expectation is at most $\bth^2(1+\bnu)\|K\Theta K\|^2\le\bth^4(1+\bnu)\,\lambda^2\|K\|^2$.
\emph{Second sum.} Expand the square into terms $\E[\theta_i\theta_k\xi_U\xi_{U'}]\,W^U_{ik}W^{U'}_{ik}$.
For $U=U'$, the Cauchy--Schwarz inequality and Lemma~\ref{lem:xi} give
$\E[\theta_i\theta_k\xi_U^2]\le\bth^2\big(\E[(1+\xi_i)^2(1+\xi_k)^2]\big)^{1/2}(\E\xi_U^4)^{1/2}
\le\bth^2(8(1+c_Q))^{1/2}c_Q^{1/2}\le3\,\bth^2(1+c_Q)$, so these terms sum to at most
$3\,\bth^2(1+c_Q)\sum_U\|W^U\|^2$. For $U\ne U'$: if neither $i$ nor $k$ lies in $U$, then $\xi_U$ is
independent of the other three factors and the expectation is zero; the same holds for $U'$. As $i$
and $k$ each lie in at most one set, the expectation can be non-zero only if one of them lies in $U$
and the other in $U'$, and then it equals $\E[(1+\xi_U)\xi_U]\,\E[(1+\xi_{U'})\xi_{U'}]=\nu_U\nu_{U'}$.
By symmetry of $W^U$ these terms sum to
$2\sum_{U\ne U'}\nu_U\nu_{U'}\sum_{i\in U,k\in U'}W^U_{ik}W^{U'}_{ik}$, which by
$|xy|\le\frac12(x^2+y^2)$ and disjointness of the sets is at most
$\bnu^2\sum_{U\ne U'}\sum_{i\in U,k\in U'}\big[(W^U_{ik})^2+(W^{U'}_{ik})^2\big]\le2\,\bnu^2\sum_U\|W^U\|^2$.
Finally $\|W^U\|^2=\tr(K_{\cdot U}K_{U\cdot}K_{\cdot U}K_{U\cdot})=\|K_{U\cdot}K_{\cdot U}\|^2=\|(K^2)_{UU}\|^2$,
so $\sum_U\|W^U\|^2\le\|K^2\|^2=\tr K^4\le\lambda^2\|K\|^2$. Collecting terms in \eqref{eq:tr4} and using
$\bth\ge1$, $\E\tr\tilde K^4\le\big[2(1+\bnu)+6(1+c_Q)+4\bnu^2\big]\bth^4\lambda^2\|K\|^2$. The last
display of the proposition uses $\|K\|^2\le\frac12V$.
\end{proof}

\subsection{A5 and the conditions of KSS}

KSS's Theorem~2 assumes their Assumption~1 and
(i) $V^{-1}\max_i\big((\tilde x_i'\beta)^2+(\check x_i'\beta)^2\big)\to0$,
(ii) $\lambda_1^2/\sum_\ell\lambda_\ell^2\to0$, where $\lambda_1,\dots,\lambda_r$ are the non-zero eigenvalues of
$\tilde A=S_{xx}^{-1/2}AS_{xx}^{-1/2}$, $|\lambda_1|$ the largest, $\tilde x_i=AS_{xx}^{-1}x_i$ and
$\check x_i=\sum_\ell M_{i\ell}\frac{B_{\ell\ell}}{1-P_{\ell\ell}}x_\ell$.

\begin{lemma}\label{lem:kss}
Let $C$ be the weights \eqref{eq:ustat}, and suppose $\underline\sigma^2\le\sigma_i^2\le\bar\sigma^2$ and
$P_{ii}\le\bar P<1$ for all $i$. Then
\[
\frac{\lambda_{\max}^2}{V}\le\frac{\bar\sigma^4}{2\,\underline\sigma^4(1-\bar P)^2}\cdot\frac{\lambda_1^2}{\sum_\ell\lambda_\ell^2},
\qquad
\max_U\phi_U\le\frac{\bar T}{V}\Big(8\,\bar\sigma^2\max_i\big[(\tilde x_i'\beta)^2+(\check x_i'\beta)^2\big]+2\lambda_{\max}^2\Big),
\]
where $\bar T$ is the largest number of rows of a stayer. So \textup{(ii)} implies condition \textup{(b)} of
Section~\ref{sec:results}, and, when $\bar T$ is bounded, \textup{(i)} and \textup{(ii)} imply condition \textup{(a)}.
\end{lemma}

\begin{proof}
Let $Z=XS_{xx}^{-1/2}$, so that $Z'Z=I$, $P=ZZ'$ and $B=Z\tilde AZ'$, where $B$ is the matrix with entries
$B_{i\ell}$. $B$ and $\tilde A$ have the same non-zero eigenvalues, so $\|B\|_\op=|\lambda_1|$ and
$\|B\|^2=\sum_\ell\lambda_\ell^2$. With $D_B=\diag(B_{ii}/M_{ii})$, $C=B-\frac12(MD_B+D_BM)$, and since $PM=0$
and $PBP=B$, $PCP=B$; a projection does not increase the Frobenius norm, so $\|C\|\ge\|B\|$. Also
$|B_{ii}|=|(Z\tilde AZ')_{ii}|\le|\lambda_1|\,(ZZ')_{ii}=|\lambda_1|\,P_{ii}$, so
$\|D_B\|_\op\le|\lambda_1|\bar P/(1-\bar P)$ and $\|C\|_\op\le\|B\|_\op+\|D_B\|_\op\le|\lambda_1|/(1-\bar P)$. Now
$K=\Sigma^{1/2}C\Sigma^{1/2}$ with $\Sigma=\diag(\sigma_i^2)$, so
$\lambda_{\max}\le\bar\sigma^2\|C\|_\op\le\bar\sigma^2|\lambda_1|/(1-\bar P)$ and
$V\ge2\|K\|^2\ge2\underline\sigma^4\|C\|^2\ge2\underline\sigma^4\sum_\ell\lambda_\ell^2$. This is the first
inequality. For the second, $M\mu=0$ gives $m=C\mu=B\mu-\frac12MD_B\mu$, and $B\mu=(\tilde x_i'\beta)_i$,
$MD_B\mu=(\check x_i'\beta)_i$, so $m_i^2\le2(\tilde x_i'\beta)^2+\frac12(\check x_i'\beta)^2$; moreover
$\sum_\ell K_{i\ell}^2=(K^2)_{ii}\le\lambda_{\max}^2$. Hence
$\phi_UV=\sum_{i\in U}\big(4\sigma_i^2m_i^2+2(K^2)_{ii}\big)$ is at most $|U|$ times the bracket.
\end{proof}

\subsection{Proposition~\ref{prop:q1} and Corollary~\ref{cor:q1}: one dominant direction}

\paragraph{Setting.} In addition to S1, S3, S4 and S5, let $\omega\in\mathbb{R}^n$ be non-random with
$\sum_i\omega_i^2=1$ (in the text $\omega_i=w_{i1}$; the symbol $w$ is in use), and let $C^{\rm r}$ be a matrix that
satisfies S2 (in the text, the weights of the remainder: they are formed as in \eqref{eq:ustat} from $A^{\rm r}$,
so Lemma~\ref{lem:stayer} applies to them). Put $\hat b=\omega'y$ and
$\hat\theta_{\rm r}=\sum_i\sum_\ell C^{\rm r}_{i\ell}y_iy_\ell$, and let $m^{\rm r}=C^{\rm r}\mu$,
$a^{\rm r}_i=\sigma_im_i^{\rm r}$, $K^{\rm r}_{i\ell}=\sigma_iC^{\rm r}_{i\ell}\sigma_\ell$,
$w^{\rm r}=a^{\rm r}+K^{\rm r}u$, $\lambda_{\rm r}=\|K^{\rm r}\|_\op$, and $k^{\rm r}_{UU}$, $k^{\rm r}_U$,
$\rho^{\rm r}_U$, $\phi^{\rm r}_U$ be the quantities of the first subsection formed with $C^{\rm r}$; in particular
$V_{\rm r}:=\V[\hat\theta_{\rm r}]=4\|a^{\rm r}\|^2+2\|K^{\rm r}\|^2>0$ and $2\rho^{\rm r}_U\le\phi^{\rm r}_UV_{\rm r}$.
Let $\varpi_i=\omega_i\sigma_i$.

\paragraph{The covariance matrix.} $\hat b-\E\hat b=\varpi'u$ and
$\hat\theta_{\rm r}-\E\hat\theta_{\rm r}=2a^{{\rm r}\prime}u+u'K^{\rm r}u$. $K^{\rm r}$ has zero diagonal, so by
Lemma~\ref{lem:forms}(b) $\E[(\varpi'u)(u'K^{\rm r}u)]=0$ with second moments only, and
\[
\Sigma_{bb}=\|\varpi\|^2=\sum_i\omega_i^2\sigma_i^2,\qquad
\Sigma_{b{\rm r}}=2\,\varpi'a^{\rm r}=2\sum_i\omega_im_i^{\rm r}\sigma_i^2,\qquad \Sigma_{\rm rr}=V_{\rm r},
\]
which is \eqref{eq:q1cov}; no third moment enters. The estimates are
$\hat\Sigma_{bb}^{\rm lo}=\sum_i\omega_i^2\hat\sigma_i^2$,
$\hat\Sigma_{b{\rm r}}=2\sum_i\omega_i\hat m_i^{\rm r}\st_i^2$ with $\hat m^{\rm r}=C^{\rm r}y$, and
$\hat\Sigma_{\rm rr}$, which is \eqref{eq:VhatA} formed with $C^{\rm r}$; $\hat\Sigma^{\rm lo}$ is the matrix with
these three entries. Let
$d_{\rm r}(\Delta)=2\sum_{i\in R}\omega_im_i^{\rm r}\Delta_i$, and let $b_{\rm r}(\Delta)$ be \eqref{eq:b} formed with
$C^{\rm r}$ and $m^{\rm r}$.

\begin{proof}[Proof of Proposition~\ref{prop:q1}]
\emph{(i) $\hat\Sigma_{bb}^{\rm lo}$.} Lemma~\ref{lem:kssvar} with $t_i=\omega_i^2$ gives
$\E\hat\Sigma_{bb}^{\rm lo}=\Sigma_{bb}$ and
$\Var(\hat\Sigma_{bb}^{\rm lo})\le c_\sigma\sum_i\omega_i^4\le c_\sigma\max_i\omega_i^2$. By K1,
$\Sigma_{bb}\ge\underline\sigma^2\sum_i\omega_i^2=\underline\sigma^2$. So
$\E[(\hat\Sigma_{bb}^{\rm lo}/\Sigma_{bb}-1)^2]\le c_b\max_i\omega_i^2$ with $c_b=c_\sigma/\underline\sigma^4$. No
model and no assumption on stayers is used.

\emph{(ii) $\hat\Sigma_{b{\rm r}}$: the mean.} Since $\sigma_i\hat m^{\rm r}_i=w^{\rm r}_i$ and
$\st_i^2=\sigma_i^2\vartheta_i(1+\xi_i)$, with $\vartheta_i=1$ where $\xi_i\ne0$,
\[
\hat\Sigma_{b{\rm r}}=Q_1+Q_2,\qquad Q_1=2\,\varpi'\Theta\,w^{\rm r},\qquad
Q_2=2\sum_{U\in\mathcal{U}}Y_U,\quad Y_U=\xi_U\,\varpi_U'w^{\rm r}_U .
\]
$\E Q_1=2\varpi'\Theta a^{\rm r}=2\sum_i\omega_im_i^{\rm r}\sigma_i^2\vartheta_i=\Sigma_{b{\rm r}}+d_{\rm r}(\Delta)$,
because $\sigma_i^2\vartheta_i=\sigma_i^2+\Delta_i$ on $R$. For $Q_2$ split $w^{\rm r}_U=g_U+h_U$ as in
\eqref{eq:gh}, with $h_U=K^{\rm r}_{UU}u_U$ and $g_U=a^{\rm r}_U+K^{\rm r}_{U,-U}u_{-U}$. Given $u_{-U}$, $g_U$ is
fixed, so by Lemmas~\ref{lem:xi} and~\ref{lem:mixed}
\begin{equation}\label{eq:Ycond}
\E\big[Y_U\,\big|\,u_{-U}\big]=\E[\xi_U]\,\varpi_U'g_U+\varpi_U'K^{\rm r}_{UU}\,\E[\xi_Uu_U]=0 .
\end{equation}
Hence $\E Q_2=0$ and $\E\hat\Sigma_{b{\rm r}}=\Sigma_{b{\rm r}}+d_{\rm r}(\Delta)$. The zero third moments of
stayers' errors are used here, through $\E[\xi_Uu_U]=0$; movers' third moments play no role, because their
rows receive non-random values. If the errors of stayer $g$ have a common third moment $\mu_{3g}$, then
$\E[\xi_Uu_j]=\E u_j^3/|U|$, and with Lemma~\ref{lem:stayer}(b) the mean of $\hat\Sigma_{b{\rm r}}$ shifts by
$2\sum_g(T_g-1)\,w_{g1}\,c_g^{\rm r}\,\mu_{3g}$.

\emph{(iii) $\hat\Sigma_{b{\rm r}}$: the variance.} $Q_1-\E Q_1=2\,(K^{\rm r}\Theta\varpi)'u$, so
$\Var(Q_1)=4\|K^{\rm r}\Theta\varpi\|^2\le4\,\bth^2\lambda_{\rm r}^2\,\Sigma_{bb}$. For $Q_2$,
$\Var(Q_2)=4\sum_U\E Y_U^2+4\sum_{U\ne U'}\E[Y_UY_{U'}]$.

\emph{Terms of the same set.} By the Cauchy--Schwarz inequality and $\|w^{\rm r}_U\|^2\le2\|g_U\|^2+2\|h_U\|^2$,
$Y_U^2\le2\,\|\varpi_U\|^2\,\xi_U^2\big(\|g_U\|^2+\|h_U\|^2\big)$. $g_U$ is independent of $\xi_U$, so
$\E[\xi_U^2\|g_U\|^2]=\nu_U\big(\|a^{\rm r}_U\|^2+\|K^{\rm r}_{U,-U}\|^2\big)$. And
$\|h_U\|^2=u_U'(K^{\rm r}_{UU})^2u_U$ with a positive semidefinite matrix of trace $k^{\rm r}_{UU}$, so by the
Cauchy--Schwarz inequality, Lemma~\ref{lem:xi} and Lemma~\ref{lem:forms}(e),
$\E[\xi_U^2\|h_U\|^2]\le(\E\xi_U^4)^{1/2}(\E\|h_U\|^4)^{1/2}\le(\E\xi_U^4)^{1/2}\big(\E(u_U'(K^{\rm r}_{UU})^2u_U)^4\big)^{1/4}
\le c_Q^{1/2}\big(8(c_Q+1)\big)^{1/4}k^{\rm r}_{UU}\le2\,(c_Q+1)\,k^{\rm r}_{UU}$. Since
$\|a^{\rm r}_U\|^2+\|K^{\rm r}_{U,-U}\|^2+k^{\rm r}_{UU}=\rho^{\rm r}_U$,
\[
\E Y_U^2\le c_Y\,\|\varpi_U\|^2\,\rho^{\rm r}_U,\qquad c_Y=2\,\bnu+4\,(c_Q+1).
\]
\emph{Terms of two different sets.} Fix $U\ne U'$, let $\Lambda=K^{\rm r}_{UU'}$ and let $O$ be the set of rows outside
$U\cup U'$; put $\bar g_U=a^{\rm r}_U+K^{\rm r}_{UO}u_O$, so that $w^{\rm r}_U=\bar g_U+h_U+\Lambda u_{U'}$, and
\[
Y_U=Y_U^0+Y_U^1,\qquad Y_U^0=\xi_U\,\varpi_U'(\bar g_U+h_U),\qquad Y_U^1=\xi_U\,\varpi_U'\Lambda\,u_{U'},
\]
and likewise for $U'$ with $\Lambda'$. $Y_U^0$ is a function of $(u_U,u_O)$ and $Y_{U'}^0$ of $(u_{U'},u_O)$. As in
\eqref{eq:Ycond}, $\E[Y_U^0\mid u_O,u_{U'}]=0$, so $\E[Y_U^0Y_{U'}^0]=0$. Next,
$Y_U^0\,Y_{U'}^1=\xi_{U'}\cdot\xi_U\,\varpi_U'(\bar g_U+h_U)\,(\Lambda\varpi_{U'})'u_U$; given $u_O$, the second
factor depends on $u_U$ only and the first on $u_{U'}$ only, the two blocks are independent, and $\E\xi_{U'}=0$: so
$\E[Y_U^0Y_{U'}^1]=0$, and by symmetry $\E[Y_U^1Y_{U'}^0]=0$. Finally
$Y_U^1Y_{U'}^1=\big(\xi_U\,u_U'\Lambda\varpi_{U'}\big)\big(\xi_{U'}\,u_{U'}'\Lambda'\varpi_U\big)$ is a product of a
function of $u_U$ and a function of $u_{U'}$, each with mean zero because $\E[\xi_Uu_U]=0$. All products are
integrable (polynomials of degree at most four in the errors of each of the two sets, and at most two in the
errors outside them). So
$\E[Y_UY_{U'}]=0$ for $U\ne U'$: the cross terms vanish exactly.

\emph{Assembly.} The sets are disjoint, so $\sum_U\|\varpi_U\|^2\le\Sigma_{bb}$, and
$\rho^{\rm r}_U\le\frac12\phi^{\rm r}_UV_{\rm r}$. Hence
\[
\Var(Q_2)=4\sum_U\E Y_U^2\le4\,c_Y\max_U\rho^{\rm r}_U\sum_U\|\varpi_U\|^2\le2\,c_Y\max_U\phi^{\rm r}_U\,\Sigma_{bb}\,\Sigma_{\rm rr},
\]
and with $\Var(Q_1+Q_2)\le2\Var Q_1+2\Var Q_2$ and $\bth\ge1$,
\[
\E\big[(\hat\Sigma_{b{\rm r}}-\Sigma_{b{\rm r}}-d_{\rm r}(\Delta))^2\big]\le
c_{b{\rm r}}\,\bth^2\Big(\max_U\phi^{\rm r}_U+\frac{\lambda_{\rm r}^2}{\Sigma_{\rm rr}}\Big)\Sigma_{bb}\,\Sigma_{\rm rr},
\qquad c_{b{\rm r}}=\max(8,\,4c_Y).
\]

\emph{(iv) $\hat\Sigma_{\rm rr}$.} $C^{\rm r}$ satisfies S2, so Propositions~\ref{prop:mean} and~\ref{prop:conc} apply to
$\hat\theta_{\rm r}$ and its variance estimate as they stand: $\E\hat\Sigma_{\rm rr}=\Sigma_{\rm rr}+b_{\rm r}(\Delta)$ and
$\E[(\hat\Sigma_{\rm rr}-\Sigma_{\rm rr}-b_{\rm r}(\Delta))^2]\le c'\bth^2(\max_U\phi^{\rm r}_U+\lambda_{\rm r}^2/\Sigma_{\rm rr})\,\Sigma_{\rm rr}^2$.

\emph{(v) The matrix.} Let $D=\diag(\Sigma_{bb}^{1/2},\Sigma_{\rm rr}^{1/2})$, $R=D^{-1}\Sigma D^{-1}$, the
correlation matrix with off-diagonal entry $\rho$, and $\hat R=D^{-1}\hat\Sigma^{\rm lo}D^{-1}$. Under the
conditions of the last part of the proposition, (i), (iii), (iv) and Chebyshev's inequality give $\hat R-R\to0$
entry by entry in probability. Now
$\Sigma^{-1/2}\hat\Sigma^{\rm lo}\Sigma^{-1/2}-I=\Sigma^{-1/2}D\,(\hat R-R)\,D\Sigma^{-1/2}$ and
$\|\Sigma^{-1/2}D\|_\op^2=\|D\Sigma^{-1}D\|_\op=\|R^{-1}\|_\op=1/(1-|\rho|)$, so
\[
\big\|\Sigma^{-1/2}\hat\Sigma^{\rm lo}\Sigma^{-1/2}-I\big\|\le\frac{\|\hat R-R\|}{1-|\rho|}\to0
\]
in probability when $|\rho|\le\bar\rho<1$.

\emph{(vi) Fitted values.} Let the movers' rows receive fitted values as in Proposition~\ref{prop:fitted}, and use
its notation: $v^*_i=\psi_i(\eta^*)$, $e_i=\hat v_i-v_i^*$, $\delta=\hat\eta-\eta^*$ and the event $\Omega$.
$\hat\Sigma_{bb}^{\rm lo}$ does not involve them. $\hat\Sigma_{\rm rr}$ is covered by Proposition~\ref{prop:fitted} applied to
$C^{\rm r}$. $\hat\Sigma_{b{\rm r}}$ is linear in the movers' values, and on $\Omega$
\[
\big|\hat\Sigma_{b{\rm r}}(\hat v)-\hat\Sigma_{b{\rm r}}(v^*)\big|=\Big|2\sum_{i\in R}\omega_i\hat m_i^{\rm r}e_i\Big|
\le2L\|\delta\|\Big(\sum_{i\in R}\omega_i^2\Big)^{1/2}\Big(\sum_{i\in R}(\hat m_i^{\rm r})^2\Big)^{1/2}
\le\frac{2L\|\delta\|}{\underline\sigma^{2}}\,\Sigma_{bb}^{1/2}\,\|w^{\rm r}_R\|,
\]
and $\E\|w^{\rm r}_R\|^2=\|a^{\rm r}_R\|^2+\|K^{\rm r}_{R\cdot}\|^2\le\frac12\Sigma_{\rm rr}$. By the Cauchy--Schwarz inequality
and $\E\|\delta\|^2\le c_\eta/n_0$, the expectation of the left-hand side over $\Omega$ is at most
$\sqrt2\,L\,c_\eta^{1/2}\underline\sigma^{-2}\,n_0^{-1/2}(\Sigma_{bb}\Sigma_{\rm rr})^{1/2}$. So the entry of $\hat R$ changes
by $O_p(n_0^{-1/2})$, and (v) applies with $\Delta_i=v_i^*-\sigma_i^2$. For the model of the text,
Corollary~\ref{cor:transfer} supplies the conditions with $n_0=n_{\cR}$.
\end{proof}

\paragraph{A sufficient condition for the remainder.} The remainder's weights are those of \eqref{eq:ustat} for
$\tilde A^{\rm r}=\tilde A-\lambda_1q_1q_1'$, whose non-zero eigenvalues are $\lambda_2,\dots,\lambda_r$. Lemma~\ref{lem:kss}
therefore applies to the remainder with $\lambda_2$ in place of $\lambda_1$, $\sum_{\ell\ge2}\lambda_\ell^2$ in place of
$\sum_\ell\lambda_\ell^2$, $\Sigma_{\rm rr}$ in place of $V$, and with $\tilde x_i$, $\check x_i$ formed from $A^{\rm r}$ and
$B^{\rm r}$ (KSS's $\tilde x_{i1}$ and $\check x_{i1}$): if the years per stayer are bounded,
$\lambda_2^2/\sum_{\ell\ge2}\lambda_\ell^2\to0$ and
$\Sigma_{\rm rr}^{-1}\max_i\big((\tilde x_{i1}'\beta)^2+(\check x_{i1}'\beta)^2\big)\to0$ imply
$\max_U\phi^{\rm r}_U+\lambda_{\rm r}^2/\Sigma_{\rm rr}\to0$. The first of these is stronger than the corresponding part of
KSS's Assumption~2, $\lambda_2^2/\sum_{\ell\ge1}\lambda_\ell^2\to0$: when $\lambda_1$ carries almost all of the sum, the
second eigenvalue can be negligible against the whole and still dominate the rest.

\paragraph{Units.} KSS's Lemma~7 asks for $\Sigma^{-1}\hat\Sigma\to I$. The interval is unchanged when $\hat b_1$ is
multiplied by a non-random number $s>0$ and $\lambda_1$ divided by $s^2$ (the ellipse and the function
$\lambda_1b^2+\theta_{\rm r}$ transform together, and $\hat\kappa$ does not change), and the conclusion of their
Theorem~3, which is about the standardised vector, is unchanged too. With $s^2=\Sigma_{\rm rr}/\Sigma_{bb}$ and
$T=\diag(s,1)$ the transformed matrices are $\Sigma^*=T\Sigma T=\Sigma_{\rm rr}R$ and
$\hat\Sigma^*=T\hat\Sigma^{\rm lo}T=\Sigma_{\rm rr}\hat R$, so $(\Sigma^*)^{-1}\hat\Sigma^*-I=R^{-1}(\hat R-R)$: for
$|\rho|\le\bar\rho<1$ their condition, in these units, is the entrywise convergence $\hat R-R\to0$ shown above.
In the original units $\Sigma^{-1}\hat\Sigma^{\rm lo}\to I$ need not hold when the two variances are of different
orders, and it is not needed.

\begin{proof}[Proof of Corollary~\ref{cor:q1}]
\emph{(a)} Write $\hat\Sigma$ for $\hat\Sigma^{\rm lo}$, positive definite with correlation $\hat\rho$,
$\hat\rho^2\le2/3$, and $\hat\Sigma^{+}$ for the same matrix with its first diagonal entry raised by
$\delta\ge0$, as \eqref{eq:q1est} does. The two intervals are ranges of the same function
$\lambda_1b^2+\theta_{\rm r}$ over two ellipses with the same centre $(\hat b_1,\hat\theta_{\rm r})$, so it suffices
that the ellipse of $\hat\Sigma^{+}$ contains that of $\hat\Sigma$; the sign of $\lambda_1$ plays no role, and for
$\lambda_1=0$ both curvatures are zero. $\hat\Sigma^{+}-\hat\Sigma$ is positive semidefinite, so
$x'(\hat\Sigma^{+})^{-1}x\le x'\hat\Sigma^{-1}x$ for every $x$. The curvature \eqref{eq:kappa} can be written
$\kappa=2|\lambda_1|\,\Sigma_{bb}^{3/2}/\det(\Sigma)^{1/2}$, and
\[
\frac{\partial\log\kappa}{\partial\Sigma_{bb}}=\frac{3}{2\,\Sigma_{bb}}-\frac{\Sigma_{\rm rr}}{2\det\Sigma}
=\frac{2\,\Sigma_{bb}\Sigma_{\rm rr}-3\,\Sigma_{b{\rm r}}^2}{2\,\Sigma_{bb}\det\Sigma},
\]
which is non-negative exactly when $\rho^2\le2/3$. Raising $\Sigma_{bb}$ lowers $\rho^2$, so the condition holds
along the way from $\hat\Sigma$ to $\hat\Sigma^{+}$, and $\hat\kappa^{+}\ge\hat\kappa$. For fixed $Z_1,Z_2$ the
quantity $(Z_1^2+(|Z_2|+r)^2)^{1/2}-r$ is non-increasing in $r=1/\kappa\ge0$ (its derivative is
$(|Z_2|+r)/(Z_1^2+(|Z_2|+r)^2)^{1/2}-1\le0$), so the critical value is non-decreasing in $\kappa$:
$z_{\alpha,\hat\kappa^{+}}\ge z_{\alpha,\hat\kappa}$. Hence the ellipse of $\hat\Sigma^{+}$ contains the ellipse of
$\hat\Sigma$. The bound is on the correlation before the entry is raised: if $\hat\rho^2>2/3$, raising the entry
can lower the curvature and with it the critical value, and an end point can move inwards.

\emph{(b)} By Proposition~\ref{prop:q1} and $|\rho|\le\bar\rho<(2/3)^{1/2}$, $\hat\Sigma^{\rm lo}$ is positive
definite with $\hat\rho_{\rm lo}^2\le2/3$ with probability tending to one. On that event the interval with
\eqref{eq:q1est} contains the interval with $\hat\Sigma^{\rm lo}$ by (a). Under the hypotheses of KSS's Theorem~3
and the consistency of $\hat\Sigma^{\rm lo}$ (in the units of the previous paragraph), their Lemma~7 gives the
latter a limiting coverage of at least $1-\alpha$. Hence the same bound holds for the former, whether or not
the first entry of \eqref{eq:q1est} is consistent.
\end{proof}

\subsection{What is not proved}

The propositions assume exact weights, a profile evaluated at fixed characteristics, known coefficients of
controls and a fixed centring constant. With projected leverages the weights have a small non-zero diagonal on
movers' rows, and no version of the propositions for these weights is proved here. Nor are: the profile evaluated
at estimated effects (the paragraph after Corollary~\ref{cor:transfer}); floors that bind in large samples; the
effects of the estimated mean and of estimated coefficients of controls on the propositions (Section~\ref{sec:controls}
gives the variance of the fixed-offset estimate, including the estimation of the control coefficients to first
order; its estimator is checked numerically in one Gaussian design only; and the additional term of the joint
estimate is not covered);
asymptotic normality of
$\hat\theta$, which is KSS's Theorem~2; and, for Section~\ref{sec:q1}, the joint normality of $\hat b_1$ and the
remainder and the coverage of the interval given a consistent covariance matrix, which are KSS's Theorem~3 and
Lemma~7. Corollary~\ref{cor:q1} needs the correlation of $\hat b_1$ and the remainder to stay below
$(2/3)^{1/2}$ in absolute value; nothing is shown for larger correlations, for a leading direction that rests on
a fixed handful of rows, or for an iteration that has not found the leading direction.
